\chapter{Competition Policy and Regulation}\label{ch:b3:competition-policy}

Competition authorities decide which mergers may go ahead and which contracts
between firms are allowed; regulators set the prices of firms that face no
competition at all. Both need tools that turn the models of the last chapters into
decisions: a way to define a market and measure its concentration, a test for
when a contract between a supplier and its retailers harms buyers, a rule for the
price of a natural monopoly. This chapter builds those tools, shows how firms respond
to the rules, and ends with the errors any enforcer makes.

\begin{openingchart}{U.S. domestic airlines: annual rates of change of the average
real yield (revenue per passenger-mile) and of traffic (revenue passenger-miles), in
three periods as the source splits them: 1966--76, before the 1978 deregulation;
1976--86, the decade around it; and the two decades that followed.}{Borenstein and Rose (2007), NBER Working Paper 13452, section on
price levels}
\begin{tikzpicture}
\begin{axis}[omrecord, ybar, bar width=9pt, width=0.6\linewidth,
  height=0.4\linewidth, xmin=-0.6, xmax=2.6, ymin=-5, ymax=10,
  xtick={0,1,2}, xticklabels from table={figdata/bachelor-3/30-competition-policy/out/yields.csv}{period},
  x tick label style={font=\scriptsize}, ylabel={Percent a year},
  legend pos=north east, legend style={font=\footnotesize}]
  \addplot[fill=omThm, draw=omThm, area legend] table[x=pos, y=yield] {figdata/bachelor-3/30-competition-policy/out/yields.csv};
  \addplot[fill=omDef!60, draw=omDef, area legend] table[x=pos, y=traffic] {figdata/bachelor-3/30-competition-policy/out/yields.csv};
  \legend{Real yield, Traffic}
\end{axis}
\end{tikzpicture}
\end{openingchart}

% ledger: b3cp:br.yield
In 1976--86, the decade around deregulation, real fares per mile fell 3.4~\% a year and traffic grew
8.2~\% a year. Such figures are often quoted as the success of deregulation. But in the
regulated decade before, real yields already fell 2.2~\% a year and traffic grew faster,
9.1~\% a year. The chart's lesson is the chapter's: judging a policy needs a
counterfactual, not a before and after.

\begin{recall}
Natural monopoly and the average-cost and price-cap rules, the Lerner index and
Cournot competition are the Year 1 volume's. Capture by regulated firms is
\cref{ch:b3:interest-groups-and-regulation}'s; mergers and their simulation are
\cref{ch:b3:oligopoly-and-collusion}'s.
\end{recall}

\begin{notation}[Symbols of this chapter]\label{not:b3:competition-policy:symbols}
$s_i$ is firm $i$'s market share, as a fraction; $\varepsilon$ the market elasticity of demand; $X$ a
price rise or a price-cap offset, as a fraction; $m$ a margin $(P - c)/P$; $s$ the allowed rate of return
on capital and $r$ its cost.
\end{notation}

\section{Markets and concentration}\label{sec:b3:competition-policy:markets}

\begin{definition}[Relevant market; hypothetical monopolist test]\label{def:b3:competition-policy:market}
\index{relevant market}\index{hypothetical monopolist test}
A \emph{relevant market} is the smallest group of products and places over which a
hypothetical monopolist would find it profitable to impose a small but significant
and non-transitory increase in price (SSNIP), typically 5~\%. The \emph{hypothetical
monopolist test} widens a candidate group until that rise pays.
\end{definition}

\begin{proposition}[Critical loss]\label{prop:b3:competition-policy:critical-loss}
A monopolist with margin $m$ that raises its price by the fraction $X$ breaks even if it loses the
share $X/(X + m)$ of its sales. The rise pays if the actual loss is smaller.
\end{proposition}

\begin{proof}
Profit per unit rises from $mP$ to $(m + X)P$. With a loss $\ell$ of sales, profit is unchanged when
$(1 - \ell)(m + X) = m$, that is $\ell = X/(X + m)$.
\end{proof}

\begin{omchart}{Critical loss of a hypothetical monopolist's price rise against its
margin, for rises of 5~\% and 10~\%.}{The formula of the proposition}
\begin{tikzpicture}
\begin{axis}[omtime, width=0.58\linewidth, height=0.36\linewidth, xmin=5, xmax=80,
  ymin=0, ymax=70, xlabel={Margin (\% of price)}, ylabel={Critical loss (\%)},
  legend pos=north east, legend style={font=\footnotesize}]
  \addplot[omDef, thick] table[x=m, y=x5] {figdata/bachelor-3/30-competition-policy/out/critical_loss.csv};
  \addplot[omThm, thick] table[x=m, y=x10] {figdata/bachelor-3/30-competition-policy/out/critical_loss.csv};
  \legend{5~\% rise, 10~\% rise}
\end{axis}
\end{tikzpicture}
\end{omchart}
% figdata: critical_loss

High margins make the critical loss small: a firm with a large margin loses much
from each lost sale. High margins also suggest that demand facing each firm is
inelastic, so the actual loss may be small too; the test needs both numbers.

\begin{measure}[The SSNIP test]
Start from the merging firms' products and region. Estimate how much of their sales a
5~\% rise would lose --- from switching surveys, diversion ratios, past responses to
price changes. Compare with the critical loss. If the loss is larger, add the closest
substitute and repeat. The group at which the rise first pays is the market.
\end{measure}

\begin{definition}[Herfindahl--Hirschman index]\label{def:b3:competition-policy:hhi}
\index{Herfindahl--Hirschman index}
The \emph{Herfindahl--Hirschman index} (HHI) is the sum of the squared market shares of
the firms in a market, with shares in percent, so that it runs from near 0 to 10{,}000.
\end{definition}

\begin{proposition}[Concentration and markups]\label{prop:b3:competition-policy:hhi-lerner}
In a Cournot equilibrium with constant marginal costs, the share-weighted average
Lerner index equals $\sum_i s_i^2/\varepsilon$: the HHI, with shares as fractions, divided by the elasticity of demand.
\end{proposition}

\begin{proof}
Firm $i$ chooses $q_i$ to maximise $P(Q)q_i - c_i q_i$: $P + P'(Q)q_i = c_i$, so $(P - c_i)/P = -P'(Q)q_i/P = s_i/\varepsilon$. Weighting by $s_i$ and summing,
$\sum_i s_i(P - c_i)/P = \sum_i s_i^2/\varepsilon$.
\end{proof}

\begin{dated}{2026}
The U.S. merger guidelines of 18 December 2023 treat a market with an HHI above 1{,}800 as
highly concentrated and an increase of more than 100 as significant. A merger that
produces a highly concentrated market with an increase above 100 is presumed to
lessen competition substantially, and so is one that creates a firm with a share
above 30~\% with an increase above 100. Five equal firms give an HHI of 2{,}000. % ledger: b3cp:doj2023
\end{dated}

\section{Vertical restraints}\label{sec:b3:competition-policy:vertical}

\begin{definition}[Vertical restraint; resale price maintenance; exclusive dealing; foreclosure]\label{def:b3:competition-policy:vertical}
\index{vertical restraint}\index{resale price maintenance}\index{exclusive dealing}\index{foreclosure}
A \emph{vertical restraint} is a contract term a supplier imposes on its buyers or
retailers. Under \emph{resale price maintenance} the supplier sets the retail price, or
a minimum for it; under \emph{exclusive dealing} a retailer may not sell rivals'
products. \emph{Foreclosure} is the denial of access to inputs or customers that rivals
need.
\end{definition}

\begin{definition}[Double marginalisation]\label{def:b3:competition-policy:double-marginalisation}
\index{double marginalisation}
\emph{Double marginalisation} is the setting of a markup by each of two firms with
market power in a chain of production, so that the final price carries two markups.
\end{definition}

\begin{proposition}[Vertical integration and the retail price]\label{prop:b3:competition-policy:double-marginalisation}
With retail demand $Q = a - P$, a supplier with cost $c$ that sells at wholesale price $w$ to a monopoly
retailer leads to the retail price $(3a + c)/4$; an integrated firm charges $(a + c)/2$, which is lower, and
earns more.
\end{proposition}

\begin{proof}
The retailer maximises $(P - w)(a - P)$: $P = (a + w)/2$ and $Q = (a - w)/2$. The supplier maximises $(w - c)(a - w)/2$: $w = (a + c)/2$, so
$P = (3a + c)/4$. The integrated firm maximises $(P - c)(a - P)$: $P = (a + c)/2 < (3a + c)/4$, since $a > c$. Its profit, $(a - c)^2/4$, exceeds the two
firms' total, $(a - c)^2/8 + (a - c)^2/16 = 3(a - c)^2/16$.
\end{proof}

\omterm{def:b3:competition-policy:vertical}{Vertical restraints} can therefore help buyers: a maximum resale price, or
integration, removes the second markup. Minimum resale prices can protect
retailers' service from free riding, or can support a retailers' \omterm{def:b3:oligopoly-and-collusion:cartel}{cartel}; \omterm{def:b3:competition-policy:vertical}{exclusive
dealing} can protect a supplier's investment in its retailers, or can foreclose
rivals. Which story holds is an empirical question for each case: the
anticompetitive version predicts higher prices and lower quantities, the efficiency
version higher prices with higher quantities.

\section{Regulating a natural monopoly}\label{sec:b3:competition-policy:regulation}

\begin{definition}[Rate-of-return regulation; Averch--Johnson effect; yardstick competition]\label{def:b3:competition-policy:regulation}
\index{rate-of-return regulation}\index{Averch--Johnson effect}\index{yardstick competition}
Under \emph{rate-of-return regulation} a firm's prices are set so that its revenue
covers its costs plus an allowed return on its capital. The \emph{Averch--Johnson
effect} is the use of more capital per unit of output than cost minimisation
requires, which such a rule induces when the allowed return exceeds the cost of
capital. Under \emph{yardstick competition} each regulated firm's allowed price
depends on the costs of comparable firms, not on its own.
\end{definition}

\begin{proposition}[Averch--Johnson]\label{prop:b3:competition-policy:aj}
A firm maximising profit $R(K, L) - wL - rK$ under the binding constraint $R(K, L) - wL \le sK$ with $s > r$ uses capital and
labour as a cost-minimising firm would if capital cost $r - \lambda(s - r)/(1 - \lambda) < r$, where $0 < \lambda < 1$ is the
constraint's multiplier: it uses more capital per unit of labour than cost
minimisation requires.
\end{proposition}

\begin{proof}
The Lagrangian is $R - wL - rK - \lambda(R - wL - sK)$. Its derivatives give $(1 - \lambda)R_L = (1 - \lambda)w$ and $(1 - \lambda)R_K = r - \lambda s$, so
$R_L = w$ and $R_K = r - \lambda(s - r)/(1 - \lambda)$. The firm equates the ratio of marginal products to $w$ over a capital
price below $r$, which with convex isoquants means a higher ratio of capital to
labour.
\end{proof}

\begin{model}{Rate-of-return regulation}
\assumes The regulator enforces the allowed return continuously; the allowed
return exceeds the cost of capital.
\predicts Over-capitalisation: costs above the minimum for the output produced.
\tested Several studies of U.S. electric utilities in the 1970s found
over-capitalisation; another view holds that regulators enforced the allowed
return only when nominal prices rose; seven Ohio electric utilities from 1965 to 1982
produced less efficiently in the years when the constraint bound more, and
regulation slowed their adoption of new technology. % ledger: b3cp:ik
\breaks When prices are reset only at intervals (regulatory lag), the firm keeps
savings between reviews and the rule works like a price cap.
\end{model}

Under a price cap the price may rise each year by inflation less an offset $X$,
whatever the firm's own costs. A firm that cuts costs keeps the saving until the
next review; under rate of return, the next price falls with its costs and the
saving goes to consumers at once, so there is little reason to make it. The weekend
problem works through both.

\section{The record of deregulation}\label{sec:b3:competition-policy:deregulation}

\begin{record}[Airline deregulation, 1978]
One counterfactual for deregulated fares is the cost-based fare formula the
regulator used just before deregulation, still updated for input costs and
productivity. In 2005 actual fares were about 30~\% below the formula's, a consumer gain of
about 28 billion dollars that year with a demand elasticity of $-1.5$; but about 26~\% of economy
passengers paid more than the formula fare for their route, and planes flew 77~\% full
against the formula's 55~\%, which changes what the formula should be. % ledger: b3cp:br.sifl
An airline's share of passengers on a route and at its end airports significantly
raises its markup, and smaller rivals on the same routes gain little from it. % ledger: b3cp:bor89
\end{record}

\section{Errors of both kinds}\label{sec:b3:competition-policy:errors}

\begin{definition}[Type I and type II errors of enforcement]\label{def:b3:competition-policy:errors}
\index{type I error of enforcement}\index{type II error of enforcement}
A \emph{type I error of enforcement} condemns conduct that does not harm buyers; a
\emph{type II error of enforcement} clears conduct that does.
\end{definition}

An authority that sees evidence $e$ and believes the conduct harmful with probability
$\pi(e)$ should condemn it when the expected cost of clearing exceeds that of condemning:
$\pi(e)C_{II} > (1 - \pi(e))C_I$, that is $\pi(e) > C_I/(C_I + C_{II})$. A rule that condemns some conduct outright (a
per se rule) is the limit in which the evidence is assumed to point one way; a
case-by-case review lowers both errors at the cost of the review. Neither error is
observed directly, which is why retrospectives such as those of
\cref{ch:b3:oligopoly-and-collusion} matter.

\section{Exercises}\label{sec:b3:competition-policy:exercises}

\begin{exercise}[$\star$]\label{exo:b3:competition-policy:1}
From the opening chart, how much faster did real yields fall in the decade after
deregulation than in the decade before? What happened to traffic growth?
\end{exercise}

\begin{exercise}[$\star$]\label{exo:b3:competition-policy:2}
Five firms have shares of 30, 25, 20, 15 and 10~\%. Compute the HHI, and the HHI after the two
smallest merge. Does the 2023 presumption apply?
\end{exercise}

\begin{exercise}[$\star$]\label{exo:b3:competition-policy:3}
With $Q = 100 - P$ and supplier cost 20, find the wholesale price, the retail price and output with a
separate retailer, and with an integrated firm.
\end{exercise}

\begin{exercise}[$\star$]\label{exo:b3:competition-policy:4}
What price and quantity data would show that minimum resale prices serve a
retailers' \omterm{def:b3:oligopoly-and-collusion:cartel}{cartel} rather than retail service?
\end{exercise}

\begin{exercise}[$\star\star$]\label{exo:b3:competition-policy:5}
Find the critical loss of a 5~\% rise for a margin of 40~\%.
\end{exercise}

\begin{exercise}[$\star\star$]\label{exo:b3:competition-policy:6}
A Cournot market has the HHI of exercise 2 before the merger and demand elasticity 2.
What is the average Lerner index?
\end{exercise}

\begin{exercise}[$\star\star$]\label{exo:b3:competition-policy:7}
A regulated firm has $Q = \sqrt{KL}$, $w = r = 10$, allowed return $s = 12$ and multiplier $\lambda = 0.5$. What capital price does it act on, and
what ratio of capital to labour does it choose, against the cost-minimising one?
\end{exercise}

\begin{exercise}[$\star\star$]\label{exo:b3:competition-policy:8}
For four invented mergers, compute the HHI before and after, the change, the merged
share, and whether the 2023 presumption applies.
\end{exercise}

\begin{exercise}[$\star\star$]\label{exo:b3:competition-policy:9}
What evidence would refute the claim that deregulation lowered fares on a
particular route?
\end{exercise}

\begin{exercise}[$\star\star\star$]\label{exo:b3:competition-policy:10}
Prove that the average Lerner index of a Cournot market is the HHI over the
elasticity, and explain why a merger that raises the HHI can still lower the price if
it lowers the merged firm's cost.
\end{exercise}

\begin{exercise}[$\star\star\star$]\label{exo:b3:competition-policy:11}
A price cap starts at 100 and rises each year by inflation, 2~\%, less $X = 3~\%$. Unit cost starts at 90 and
falls 4~\% a year; output is 1{,}000 a year. Compute price and profit for five years, and compare
with a price reset each year to cost plus 10.
\end{exercise}

\begin{exercise}[$\star\star\star$]\label{exo:b3:competition-policy:12}
Wrongly condemning a contract costs 1 and wrongly clearing it costs 3. Above what
probability of harm should the authority condemn? If 20~\% of such contracts are
harmful and the evidence doubles the odds of harm, what should it do?
\end{exercise}

\section{Problem: Brenland's Water Company}\label{sec:b3:competition-policy:problem}

\begin{problem}[{Weekend problem --- Brenland's water company}]\label{pb:b3:competition-policy:water}
One firm supplies Brenland's capital with water. Daily demand is $Q = 100 - P$, in thousand
cubic metres, with $P$ in dalens per cubic metre. Output is $Q = \sqrt{KL}$; capital and labour each cost 10
a unit, so marginal cost is 20; the network's fixed cost is 600 (thousand dalens a day). % figdata: water_stats

\medskip\noindent\textbf{Part I --- no regulation.}
\begin{enumerate}
\item Show that marginal cost is 20, by minimising the cost of producing $Q$.
\item Why is this a natural monopoly?
\item Find the monopoly price, output and profit.
\item Find consumer surplus.
\item Find the deadweight loss against marginal-cost pricing.
\end{enumerate}

\medskip\noindent\textbf{Part II --- marginal and average cost.}
\begin{enumerate}[resume]
\item With price at marginal cost, find output, consumer surplus and the firm's loss.
\item Find the average-cost price and output.
\item Find consumer surplus and the deadweight loss at the average-cost price.
\item Who pays the fixed cost under each rule?
\end{enumerate}

\medskip\noindent\textbf{Part III --- rate of return.} The regulator requires 60 a day at the price
consumers pay for it, and allows revenue to cover labour and fixed costs plus 12 per unit
of capital.
\begin{enumerate}[resume]
\item Find the price, and the cost-minimising inputs and cost for 60.
\item What profit would the cost-minimising plan earn, and what does the regulator
allow it?
\item Show that the firm's allowed profit is $2K$.
\item Find the capital, labour and ratio of capital to labour the firm chooses.
\item Find its cost, and the excess over the minimum.
\item Why would the regulator have trouble seeing the excess?
\end{enumerate}

\medskip\noindent\textbf{Part IV --- a price cap.} Instead, the price is capped at the
average-cost price less $X$; the firm can cut its marginal cost from 20 to 16 by an effort that
costs 150 a day.
\begin{enumerate}[resume]
\item With $X = 0$, find the firm's profit if it makes the effort. Does it?
\item Under cost-based pricing, the price would fall to average cost with the new
marginal cost. Find it, and the firm's profit after the effort. Does it make it?
\item Find the $X$ at which consumers and the firm gain equally from the effort.
\item Find the price, consumer surplus and profit at that $X$.
\item Why must $X$ be set before the effort, and kept?
\item What happens at the next review if the regulator resets the price to the new
average cost?
\item How would \omterm{def:b3:competition-policy:regulation}{yardstick competition} help here?
\item Compare consumer surplus across monopoly, average cost and the cap.
\item Which regime gives the highest consumer surplus, and why is it not chosen?
\item State the result: consumer surplus and the firm's profit under each regime, and
the $X$ that shares the cost saving equally.
\end{enumerate}
\end{problem}
